% !TEX program = XeLaTeX
% Домашнее задание по курсу «Машинное обучение в экономике», 2025-2026.
% Компилируется XeLaTeX (Overleaf: Menu -> Compiler -> XeLaTeX). Рисунки в папке figures/.
% Шрифт Times New Roman. Если он недоступен в системе сборки (например, на Overleaf),
% замените "Times New Roman" ниже на "TeX Gyre Termes" или "Liberation Serif".
% ВАЖНО: согласно требованиям, PDF не должен содержать кода. Код -- в отдельном файле hw.py.
\documentclass[12pt,a4paper]{article}
\usepackage{fontspec}
\setmainfont{Times New Roman}[Ligatures=TeX]
\newfontfamily\cyrillicfont{Times New Roman}[Ligatures=TeX]
\usepackage{polyglossia}
\setdefaultlanguage{russian}
\setotherlanguage{english}
\usepackage{amsmath,amssymb}
\usepackage{geometry}
\geometry{left=2.5cm,right=2cm,top=2.5cm,bottom=2.5cm}
\usepackage{graphicx}
\usepackage{booktabs}
\usepackage{array}
\usepackage{caption}
\usepackage{float}
\usepackage{setspace}
\onehalfspacing
\usepackage[hidelinks]{hyperref}

\newcommand{\zadanie}[2]{\subsection*{Задание #1}\begingroup\itshape #2\par\endgroup\medskip}
\newcommand{\otvet}{\par\noindent\textbf{Ответ}\par\medskip}
\long\def\refitem#1{\par\noindent\hangindent=1.27cm #1\par\medskip}

\begin{document}

\begin{titlepage}
\centering
\vspace*{1.5cm}
{\large Национальный исследовательский университет\\ «Высшая школа экономики»}\\[0.3cm]
{\large Факультет экономических наук}\\[3cm]
{\Large\bfseries Домашнее задание по курсу}\\[0.3cm]
{\Large\bfseries «Машинное обучение в экономике»}\\[1cm]
{\large\itshape Влияние подключения малого бизнеса к маркетплейсу\\ на месячную выручку фирмы}\\[3cm]
\begin{flushright}\large
\textit{Выполнили студенты группы МЭКАН25АДЭ1:}\\[0.2cm]
Романов Алексей\\ Галкин Евгений\\[1cm]
\textit{Семинаристы:}\\ Потанин Б. С.\\ Долгих С. И.\\
\end{flushright}
\vfill {\large Москва, 2026}
\end{titlepage}

\tableofcontents
\newpage

\section{Обоснование темы}

\zadanie{1.1}{Придумайте непрерывную зависимую (целевую) переменную и эндогенную бинарную переменную воздействия. Обоснуйте, почему переменная воздействия может быть эндогенной, и создайте хотя бы одну ненаблюдаемую переменную, которая коррелирует и с целевой переменной, и с переменной воздействия.}
\otvet
Объектом исследования служат небольшие фирмы (индивидуальные предприниматели и микропредприятия), торгующие потребительскими товарами. Зависимая переменная это месячная выручка фирмы через шесть месяцев после момента возможного подключения к маркетплейсу, в тысячах рублей; обозначим ее $Y$. Это непрерывная положительная величина, что соответствует требованию к целевой переменной.

Переменная воздействия это факт подключения фирмы к маркетплейсу: $D=1$, если фирма разместила товары на площадке и начала продавать через нее, и $D=0$ иначе. Нас интересует, насколько выход на маркетплейс меняет выручку типичной небольшой фирмы.

Переменная $D$ с высокой вероятностью эндогенна. Решение о подключении принимает собственник, и оно не случайно: к маркетплейсу быстрее подключаются предприниматели с более высоким качеством менеджмента, лучшим пониманием спроса и большей готовностью осваивать новые каналы продаж. Эти же качества повышают выручку и без маркетплейса, через грамотное ценообразование, ассортиментную политику и работу с издержками. Поэтому простое сравнение выручки подключившихся и не подключившихся фирм смешивает эффект площадки с эффектом ненаблюдаемых способностей собственника. Чтобы формализовать этот механизм, мы вводим ненаблюдаемую переменную $U$, отражающую качество менеджмента и предпринимательские способности собственника. По построению $U$ положительно связана и с вероятностью подключения, и с выручкой, то есть является общей причиной $D$ и $Y$. В соответствии с заданием $U$ используется при генерации данных, но не включается в число признаков ни в одну из моделей анализа.

\zadanie{1.2}{Опишите, для чего может быть полезно изучение влияния переменной воздействия на зависимую переменную, в том числе для бизнеса и государственных органов.}
\otvet
Оценка причинного эффекта подключения к маркетплейсу полезна сразу нескольким агентам. Для самой фирмы знание величины эффекта помогает решить, выходить ли на площадку: подключение сопряжено с издержками (комиссия, логистика, упаковка, фотоконтент, продвижение карточек), и если ожидаемый прирост выручки невелик или сосредоточен лишь у части фирм, ресурсы разумнее направить в другие каналы. Для оператора маркетплейса оценка эффекта важна при настройке программ привлечения продавцов: площадка несет затраты на онбординг и временные льготы, поэтому ей выгодно приглашать в первую очередь фирмы с наибольшим ожидаемым приростом оборота. Для государственных органов результаты полезны при разработке мер поддержки малого и среднего предпринимательства и программ цифровизации: если выход на маркетплейс действительно повышает выручку, субсидирование подключения и обучение продавцов оправданы, а если эффект мал или достается и без того успешным фирмам, средства разумнее перенаправить. Корректная причинная оценка помогает избежать как недооценки, так и переоценки отдачи от таких программ.

\zadanie{1.3}{Обоснуйте наличие причинно-следственной связи между зависимой переменной и переменной воздействия. Опишите результаты не менее 2 предшествующих исследований, критически оцените их методологию и сравните с применяемыми вами методами.}
\otvet
Содержательно связь между подключением к маркетплейсу и выручкой опирается на несколько механизмов: площадка расширяет доступную аудиторию покупателей, снижает издержки поиска и выхода на новые географические рынки и предоставляет готовую инфраструктуру оплаты и доставки. Все это увеличивает число заказов и выручку.

Lendle и др. (2016) на данных eBay по 61 стране показывают, что влияние расстояния на онлайн-торговлю в среднем на 65\% слабее, чем на обычную международную торговлю сопоставимыми товарами, и связывают это с сокращением издержек поиска. Для нашей задачи это означает, что площадка снимает часть барьеров выхода на рынок и расширяет круг покупателей, что является прямым каналом роста продаж. Методологически работа опирается на гравитационную модель на агрегированных данных: ее сильная сторона это широкий межстрановой охват, однако подход наблюдательный и не позволяет оценить причинный эффект для отдельной фирмы, поскольку сопоставляются товарные потоки, а не выручка продавцов, и не контролируется отбор того, какие товары выходят в онлайн.

Couture и др. (2021) изучают общенациональную программу развития сельской электронной коммерции в Китае при помощи рандомизированного эксперимента и находят ограниченный прирост доходов сельских производителей: основной выигрыш приходится на снижение стоимости жизни. Эта работа важна как осторожный контраргумент: эффект электронной коммерции может быть неоднородным и не сводиться к автоматическому росту доходов продавцов. Рандомизация обеспечивает высокую внутреннюю валидность, но воздействием здесь выступает появление в деревне пункта электронной коммерции, а не подключение конкретной фирмы, а контекст сельского Китая ограничивает перенос результатов.

В нашем исследовании единицей анализа является фирма, воздействием служит ее подключение, а зависимой переменной ее выручка. Для идентификации эффекта мы используем инструментальную переменную и методы причинного машинного обучения (двойное машинное обучение и метаобучающие алгоритмы для неоднородных эффектов), которые подробно вводятся в разделах 4 и 5. Преимущество подхода в том, что он позволяет оценивать средний, локальный и условные средние эффекты, гибко учитывая нелинейные связи и не требуя корректной параметрической спецификации. Главное ограничение нашей работы по сравнению с цитируемыми это использование симулированных, а не реальных данных: это дает полный контроль над истинными значениями эффектов и возможность проверить качество методов, но не позволяет судить о реальной величине эффекта подключения.

\zadanie{1.4}{Придумайте хотя бы 3 контрольные переменные (минимум одна бинарная и одна непрерывная) и обоснуйте выбор каждой.}
\otvet
Мы вводим четыре контрольные переменные. Возраст фирмы $\textit{age}$ (непрерывная, годы): более старые фирмы наработали клиентскую базу и репутацию, что повышает базовую выручку, но могут быть подвержены инерции и реже осваивают новый канал. Число занятых $\textit{size}$ (непрерывная): размер отражает производственные и операционные возможности, прямо увеличивает выручку и облегчает работу с площадкой. Наличие онлайн-канала $\textit{online}$ (бинарная) равно единице, если до периода фирма уже продавала онлайн (сайт или соцсети): это повышает базовую выручку и облегчает подключение, но снижает предельный эффект площадки, так как часть онлайн-спроса уже охвачена. Расположение в крупном городе $\textit{city}$ (бинарная) это прокси доступности логистики и платежеспособного спроса: в крупных городах базовая выручка выше, но дополнительный эффект маркетплейса может быть меньше.

\zadanie{1.5}{Придумайте бинарную инструментальную переменную и обоснуйте, почему она удовлетворяет необходимым условиям.}
\otvet
Инструмент $Z$ это факт попадания фирмы в случайную волну приглашений на ускоренное подключение: $Z=1$, если фирма получила приглашение, и $Z=0$ иначе. Такие волны практикуются площадками, рассылающими части продавцов предложения о льготном онбординге (упрощенная регистрация, сниженная комиссия, помощь менеджера); состав волны формируется случайно. Инструмент удовлетворяет трем условиям. Релевантность: приглашение снижает издержки выхода и повышает вероятность подключения, то есть $\mathrm{Cov}(Z,D)\neq 0$ (проверяется на данных). Независимость: приглашения рассылаются случайно и не зависят от ненаблюдаемого качества менеджмента $U$. Исключение: само приглашение, если фирма им не воспользовалась, не меняет продажи, поэтому единственный канал влияния $Z$ на $Y$ проходит через $D$. Дополнительно выполняется монотонность: приглашение может только повысить вероятность подключения. При этих условиях инструмент идентифицирует локальный средний эффект для фирм-комплаеров, чье решение меняется под действием приглашения (Imbens, Angrist, 1994).

\zadanie{1.6}{Приведите дополнительные комментарии о целях, задачах, методологии и вкладе исследования.}
\otvet
Цель работы это на контролируемых синтетических данных воспроизвести типичную ситуацию с эндогенной переменной воздействия и сравнить, насколько разные классы методов справляются с оценкой причинного эффекта. Поскольку процесс генерации данных задан нами, истинные значения среднего, условных средних и локального среднего эффектов известны, что позволяет оценивать не только качество прогнозов, но и точность восстановления эффектов. Задачи включают построение достаточно сложного процесса генерации данных с эндогенностью, обучение и сравнение моделей классификации и регрессии, а также оценку эффектов воздействия разными методами с сопоставлением их между собой и с истинными значениями. Современная литература предлагает широкий арсенал методов оценки эффектов: от разности разностей (Котырло, 2024) до параметрических моделей с эндогенным переключением и неслучайным отбором (Коссова, Потанин, 2018, 2022) и подходов на основе мэтчинга (Коссова, Косорукова, 2023); мы сопоставляем ряд таких методов, уделяя внимание тому, как качество вспомогательных моделей машинного обучения влияет на итоговые оценки. Близкие по духу эмпирические исследования факторов эффективности российских фирм проводятся, в частности, в работе Долгих и Потанин (2023).

\section{Генерация и предварительная обработка данных}

\zadanie{2.1}{Опишите математически предполагаемый процесс генерации данных: распределения экзогенных переменных; потенциальные исходы воздействия в зависимости от инструмента; потенциальные исходы зависимой переменной; обоснуйте направления связей. Не рекомендуется использовать простые линейные модели.}
\otvet
Опишем процесс генерации данных для $i=1,\dots,n$ независимых фирм.

\textbf{Экзогенные переменные.} Ненаблюдаемое качество менеджмента $U_i\sim\mathcal{N}(0,1)$. Возраст фирмы логнормален, $\textit{age}_i=\exp(\eta_i)$, $\eta_i\sim\mathcal{N}(1.6,0.5^2)$; число занятых $\textit{size}_i=\lceil\exp(\zeta_i)\rceil$, $\zeta_i\sim\mathcal{N}(1.2,0.6^2)$. Наличие онлайн-канала зависит от качества менеджмента: $\Pr(\textit{online}_i=1\mid U_i)=\Lambda(-0.2+0.8\,U_i)$, где $\Lambda(t)=1/(1+e^{-t})$. Расположение в крупном городе $\textit{city}_i\sim\mathrm{Bernoulli}(0.5)$. Инструмент генерируется независимо от всех факторов: $Z_i\sim\mathrm{Bernoulli}(0.5)$, $Z_i\perp(U_i,\textit{age}_i,\textit{size}_i,\textit{online}_i,\textit{city}_i)$, что отражает случайный характер рассылки.

\textbf{Потенциальные исходы воздействия.} Латентный индекс склонности к подключению
\[
D^{*}_i(z)=\gamma_0+\gamma_z z+\gamma_u U_i+\gamma_{u2}U_i^2+g(\cdot)+\varepsilon^{D}_i,\quad \varepsilon^{D}_i\sim\mathcal{N}(0,1),
\]
\[
g(\cdot)=\delta_1\ln(\textit{age}_i)+\delta_2\sqrt{\textit{size}_i}+\delta_3\textit{online}_i+\delta_4\textit{city}_i+\delta_5\,\textit{online}_i\textit{city}_i.
\]
Потенциальные исходы $D_i(0)=\mathbf{1}\{D^{*}_i(0)\ge0\}$, $D_i(1)=\mathbf{1}\{D^{*}_i(1)\ge0\}$. Так как $\gamma_z>0$, для любой фирмы $D_i(1)\ge D_i(0)$ (монотонность), что разбивает совокупность на never-takers, always-takers и комплаеров. Наблюдаемое значение $D_i=(1-Z_i)D_i(0)+Z_iD_i(1)$. Качество менеджмента и наличие онлайн-канала повышают склонность к подключению ($\gamma_u>0$, $\delta_3>0$), как и размер ($\delta_2>0$); возраст связан отрицательно из-за инерции ($\delta_1<0$); квадратичный член вводит нелинейность.

\textbf{Потенциальные исходы выручки.} С гетерогенным эффектом и нелинейной базой:
\[
Y_i(d)=\mu(\cdot)+d\cdot\tau_i+\sigma\xi_i,\quad \xi_i\sim\mathcal{N}(0,1),
\]
\[
\mu(\cdot)=\beta_0+\beta_1\sqrt{\textit{size}_i}+\beta_2\ln(\textit{age}_i)+\beta_3\textit{online}_i+\beta_4\textit{city}_i+\beta_5U_i+\beta_6U_i^2,
\]
\[
\tau_i=\tau_0+\tau_1U_i+\tau_2\textit{online}_i+\tau_3\textit{city}_i+\tau_4\ln(\textit{age}_i).
\]
Наблюдаемая выручка $Y_i=(1-D_i)Y_i(0)+D_iY_i(1)$. Размер, возраст, онлайн-канал и крупный город повышают базовую выручку ($\beta_1,\dots,\beta_4>0$); качество менеджмента входит положительно ($\beta_5>0$) с нелинейностью ($\beta_6$). Поскольку $U$ входит и в уравнение подключения, и в уравнение выручки, именно она порождает эндогенность. Эффект подключения в среднем положителен ($\tau_0>0$) и неоднороден: выше у фирм с лучшим менеджментом ($\tau_1>0$), ниже у фирм с онлайн-каналом ($\tau_2<0$) и в крупных городах ($\tau_3<0$). Зависимость $\tau$ от $U$ приводит к различию локального среднего эффекта для комплаеров и среднего эффекта по совокупности (см. раздел 5). Числовые значения параметров заданы в файле \texttt{hw.py}.

\zadanie{2.2}{Симулируйте данные, приведите корреляционную матрицу и описательные статистики. Не менее 1000 наблюдений; доля единиц каждой бинарной переменной от 0.1 до 0.9.}
\otvet
Сгенерировано $n=10000$ наблюдений (объем выбран с запасом: согласно примечанию к заданию, для точной оценки эффектов воздействия по наблюдаемым значениям полезна выборка порядка десяти тысяч). Описательные статистики приведены в таблицах~\ref{tab:desc_cont} и~\ref{tab:desc_bin}: доли единиц всех бинарных переменных лежат в требуемом диапазоне $[0.1,0.9]$. Корреляционная матрица наблюдаемых переменных показана на рис.~\ref{fig:corr}.

\begin{table}[H]\centering
\caption{Описательные статистики непрерывных переменных}\label{tab:desc_cont}
\begin{tabular}{lrrrrr}\toprule
Переменная & Среднее & Ст.\,откл. & Медиана & Мин & Макс\\\midrule
$Y$ (выручка, тыс. руб.) & 200.22 & 55.80 & 196.61 & 27.44 & 426.75\\
$\textit{age}$ (возраст, лет) & 5.61 & 3.01 & 4.95 & 0.69 & 36.48\\
$\textit{size}$ (занятые) & 4.49 & 2.68 & 4.00 & 1.00 & 41.00\\\bottomrule
\end{tabular}
\end{table}

\begin{table}[H]\centering
\caption{Описательные статистики бинарных переменных}\label{tab:desc_bin}
\begin{tabular}{lrr}\toprule
Переменная & Доля единиц & Кол-во единиц\\\midrule
$D$ (подключение) & 0.488 & 4884\\
$Z$ (приглашение) & 0.495 & 4948\\
$\textit{online}$ (онлайн-канал) & 0.450 & 4497\\
$\textit{city}$ (крупный город) & 0.491 & 4914\\\bottomrule
\end{tabular}
\end{table}

\begin{figure}[H]\centering
\includegraphics[width=0.7\textwidth]{figures/corr.png}
\caption{Корреляционная матрица наблюдаемых переменных}\label{fig:corr}
\end{figure}

Структура корреляций согласуется с замыслом. Выручка $Y$ заметно связана с подключением $D$ ($0.56$), а также с размером ($0.42$) и наличием онлайн-канала ($0.44$). Инструмент $Z$ сильно коррелирует с подключением $D$ ($0.26$), но почти не связан с выручкой напрямую ($0.05$): это эмпирическое отражение условия исключения, поскольку влияние $Z$ на $Y$ проходит только через $D$. Доли типов фирм по реакции на инструмент: never-takers $0.381$, always-takers $0.355$, комплаеры $0.264$.

\zadanie{2.3}{Разделите выборку на обучающую и тестовую (тест 20--30\%).}
\otvet
Выборка разбита на обучающую (7500 наблюдений, 75\%) и тестовую (2500 наблюдений, 25\%) со стратификацией по переменной воздействия $D$, чтобы сохранить ее долю в обеих частях. Это разбиение используется в разделах 3 и 4; в разделе 5 части объединяются.

\zadanie{2.4}{Дополнительный анализ и комментарии о генерации данных.}
\otvet
Отметим, что предсказуемость подключения $D$ по наблюдаемым признакам заведомо ограничена: значимая часть вариации $D$ определяется ненаблюдаемым качеством менеджмента $U$ и случайной компонентой, которые недоступны моделям. Это сознательная особенность процесса: она делает задачу классификации в разделе 3 содержательно трудной и одновременно служит источником эндогенности, с которой работает раздел 5.

\section{Классификация}

\zadanie{3.1}{Отберите признаки для прогнозирования переменной воздействия (без целевой переменной).}
\otvet
Для прогноза подключения $D$ используются признаки $\{Z,\textit{age},\textit{size},\textit{online},\textit{city}\}$. Выручка $Y$ как целевая переменная исключена, ненаблюдаемое качество менеджмента $U$ недоступно по построению. Инструмент $Z$ включен как законный предиктор подключения. Использованы пять методов: логистическая регрессия, метод ближайших соседей, наивный байесовский классификатор, случайный лес и градиентный бустинг.

\zadanie{3.2}{Для каждого метода оцените accuracy на обучающей, тестовой выборках и по кросс-валидации; подберите гиперпараметры; сведите результаты в таблицу; обсудите переобучение.}
\otvet
Результаты приведены в таблице~\ref{tab:clf_tuning}. Точность умеренная (около $0.66$--$0.67$), что ожидаемо: значимая часть вариации $D$ определяется ненаблюдаемым $U$, недоступным моделям, поэтому верхняя граница предсказуемости ограничена. Ярче всего переобучение видно у случайного леса с параметрами по умолчанию: точность на обучении $1.000$ при тесте $0.589$. Кросс-валидационный подбор гиперпараметров (ограничение глубины и минимального размера листа) устраняет разрыв и поднимает тестовую точность до $0.660$. У логистической регрессии и наивного Байеса разрыв между обучением и тестом мал, признаков переобучения нет.

\begin{table}[H]\centering\small
\caption{Точность классификаторов до и после настройки гиперпараметров}\label{tab:clf_tuning}
\begin{tabular}{lccccc}\toprule
Метод & Acc train & Acc CV & Acc test & Лучшие гиперпарам. & Acc test (тюн.)\\\midrule
LogReg & 0.665 & 0.666 & 0.672 & $C=1$ & 0.672\\
kNN & 0.755 & 0.625 & 0.628 & $k=51$ & 0.662\\
NaiveBayes & 0.650 & 0.649 & 0.659 & var\_sm.$=10^{-9}$ & 0.659\\
RandomForest & 1.000 & 0.596 & 0.589 & depth 12, leaf 10 & 0.660\\
GradBoost & 0.688 & 0.660 & 0.674 & lr 0.1, depth 2, 200 дер. & 0.666\\\bottomrule
\end{tabular}
\end{table}

\textit{Повышенная сложность (OOB).} Для случайного леса гиперпараметры подобраны также по out-of-bag ошибке. Максимум OOB-точности достигается при глубине 6 и минимальном листе 1 (OOB $0.663$, тест $0.676$), тогда как кросс-валидация выбрала более сильную регуляризацию (глубина 12, лист 10). OOB-оценка получается почти бесплатно в ходе обучения леса и не требует отдельных разбиений, что особенно удобно при больших объемах данных; ее недостаток в том, что она специфична для бэггинга, использует только бутстреп-выборки и при сильной корреляции деревьев может быть смещена, тогда как кросс-валидация применима к любому методу и дает более универсальную оценку.

\zadanie{3.3}{Повторите подбор гиперпараметров по альтернативному критерию качества; обсудите его преимущества и недостатки.}
\otvet
В качестве альтернативного критерия использован AUC (таблица~\ref{tab:clf_auc_tuning}). В отличие от accuracy, AUC не зависит от порога классификации и оценивает качество ранжирования, что важно в нашей прикладной задаче, где порог в дальнейшем подбирается из экономических соображений. Подбор по AUC дал близкие к accuracy-настройке гиперпараметры; для логистической регрессии оптимальным оказалось чуть более сильное сжатие ($C=0.1$).

\begin{table}[H]\centering\small
\caption{Подбор гиперпараметров по AUC}\label{tab:clf_auc_tuning}
\begin{tabular}{lccc}\toprule
Метод & Лучшие гиперпарам. (AUC) & CV AUC & AUC test\\\midrule
LogReg & $C=0.1$ & 0.734 & 0.744\\
kNN & $k=51$ & 0.726 & 0.729\\
NaiveBayes & var\_sm.$=10^{-9}$ & 0.721 & 0.731\\
RandomForest & depth 12, leaf 10 & 0.717 & 0.721\\
GradBoost & lr 0.1, depth 2, 200 дер. & 0.725 & 0.739\\\bottomrule
\end{tabular}
\end{table}

\textit{Повышенная сложность (собственный критерий).} Дополнительно запрограммирован критерий ожидаемой прибыли таргетинга (см. задание 3.7): для каждого набора гиперпараметров вычисляется максимальная по порогу прибыль на кросс-валидации. При настройке случайного леса по accuracy выбирается сильная регуляризация (минимальный лист 10), а по критерию прибыли менее регуляризованная модель (минимальный лист 1): прибыль чувствительна к качеству ранжирования верхних по вероятности фирм, поэтому допускает более гибкую модель. Такой критерий уместен, когда у ошибок разная экономическая цена, и менее полезен как универсальная мера качества, поскольку зависит от заданных цен прогнозов.

\zadanie{3.4}{Постройте ROC-кривые и сравните модели по AUC; обсудите преимущества и недостатки AUC по сравнению с accuracy.}
\otvet
ROC-кривые на тестовой выборке показаны на рис.~\ref{fig:roc_clf}, значения AUC в таблице~\ref{tab:clf_auc}. Лучшее качество ранжирования у логистической регрессии (AUC $0.744$) и градиентного бустинга/CatBoost ($0.739$). AUC предпочтительнее accuracy в нашей задаче, поскольку не зависит от порога и устойчив к дисбалансу классов; его недостаток в том, что он оценивает только упорядочивание и не отражает качество калибровки вероятностей и реальные издержки конкретного порога, которые важны при принятии решений.

\begin{table}[H]\centering\small
\caption{AUC классификаторов на тестовой выборке}\label{tab:clf_auc}
\begin{tabular}{lcccccc}\toprule
& LogReg & kNN & NaiveBayes & RandomForest & GradBoost & CatBoost\\\midrule
AUC test & 0.744 & 0.729 & 0.731 & 0.721 & 0.739 & 0.739\\\bottomrule
\end{tabular}
\end{table}

\begin{figure}[H]\centering
\includegraphics[width=0.66\textwidth]{figures/roc_clf.png}
\caption{ROC-кривые классификаторов на тестовой выборке}\label{fig:roc_clf}
\end{figure}

\zadanie{3.5}{Постройте матрицу ошибок при пороге 0.5 и обсудите, какие ошибки встречаются чаще.}
\otvet
Для лучшей по AUC модели (логистическая регрессия) при пороге $0.5$ матрица ошибок на тесте приведена в таблице~\ref{tab:confusion}. Ошибки обоих типов сопоставимы по частоте: $401$ ложноположительных и $418$ ложноотрицательных. Близость их числа объясняется тем, что классы примерно сбалансированы, а порог $0.5$ нейтрален относительно типов ошибок.

\begin{table}[H]\centering
\caption{Матрица ошибок (LogReg, порог 0.5, тест)}\label{tab:confusion}
\begin{tabular}{lcc}\toprule
& Прогноз 0 & Прогноз 1\\\midrule
Факт 0 & 878 & 401\\
Факт 1 & 418 & 803\\\bottomrule
\end{tabular}
\end{table}

\zadanie{3.6}{Рассмотрите не менее двух порогов, отличных от 0.5; рассчитайте ACC, Precision, Recall; обсудите влияние порога и предпочтительный порог.}
\otvet
Таблица~\ref{tab:thresholds} показывает, что снижение порога до $0.3$ резко повышает полноту (Recall $0.904$) ценой точности (Precision $0.586$), а повышение до $0.7$ дает обратный эффект (Precision $0.783$, Recall $0.343$). Если заказчик это маркетплейс, рассылающий приглашения, более существенны ложноотрицательные прогнозы (упущенные перспективные продавцы), поэтому при равной цене контакта предпочтителен более низкий порог $0.3$, обеспечивающий высокую полноту; окончательный выбор порога делается ниже из критерия прибыли.

\begin{table}[H]\centering
\caption{Метрики при разных порогах (LogReg, тест)}\label{tab:thresholds}
\begin{tabular}{lccc}\toprule
Порог & ACC & Precision & Recall\\\midrule
0.3 & 0.641 & 0.586 & 0.904\\
0.5 & 0.672 & 0.667 & 0.658\\
0.7 & 0.633 & 0.783 & 0.343\\\bottomrule
\end{tabular}
\end{table}

\zadanie{3.7}{Опишите прикладную задачу и заказчика, задайте цены прогнозов и подберите порог, максимизирующий прибыль; сравните прибыли на тесте.}
\otvet
Заказчик это оператор маркетплейса, рассылающий приглашения на льготное подключение. Прогноз положителен, если фирма, по оценке модели, склонна подключиться. На таких фирм оператор тратит стоимость контакта $C=3$ тыс. руб.; если приглашенная фирма действительно подключается, оператор получает выгоду $G=10$ тыс. руб. Тогда истинноположительный прогноз приносит $G-C=7$, ложноположительный $-C=-3$, остальные $0$. Оптимальный порог подбирается на обучающей выборке, прибыль считается на тесте (таблица~\ref{tab:profit}). Оптимальные пороги лежат заметно ниже $0.5$ (около $0.28$--$0.38$): поскольку выгода от привлеченной фирмы превышает стоимость контакта, выгоднее приглашать шире, мирясь с частью ложноположительных. Это согласуется с экономической интуицией: при дешевом контакте и дорогом привлечении оптимально снижать порог. Наибольшую прибыль на тесте дают логистическая регрессия и наивный Байес.

\begin{table}[H]\centering\small
\caption{Оптимальные пороги и прибыль на тесте}\label{tab:profit}
\begin{tabular}{lcccc}\toprule
Метод & Порог (лин.) & Прибыль (лин.) & Порог (нелин.) & Прибыль (нелин.)\\\midrule
LogReg & 0.305 & 5368 & 0.455 & 3064\\
kNN & 0.275 & 5198 & 0.435 & 3032\\
NaiveBayes & 0.285 & 5309 & 0.445 & 2957\\
RandomForest & 0.325 & 5193 & 0.480 & 3037\\
GradBoost & 0.380 & 5235 & 0.430 & 3060\\\bottomrule
\end{tabular}
\end{table}

\textit{Повышенная сложность (нелинейная функция прибыли).} Рассмотрена также нелинейная функция, в которой ценность привлеченной фирмы вогнуто зависит от ее прогнозной вероятности подключения, $G\sum_{i:\,\hat p_i\ge t,\,D_i=1}\sqrt{\hat p_i}-C\cdot\#\{\hat p_i\ge t\}$. Она отражает убывающую отдачу и сдвигает оптимальный порог вверх (около $0.43$--$0.48$), делая таргетинг более консервативным.

\zadanie{3.8}{Опишите связи переменных в виде DAG, обучите структуру байесовской сети и сравните точность вашего и обученного DAG.}
\otvet
Наш экспертный DAG задает направленные связи от признаков $\{Z,\textit{online},\textit{city},\textit{age},\textit{size}\}$ к подключению $D$ (непрерывные признаки дискретизованы по терцилям). Структура, обученная алгоритмом восхождения по холмам с критерием BIC, выделила содержательно осмысленные ребра $Z\to D$ и $\textit{online}\to D$ (а также связи $D$ с размером и возрастом). Точность прогноза $D$ на тесте у обеих сетей совпала и составила $0.669$ (таблица~\ref{tab:bayesnet}), что близко к лучшим дискриминативным моделям. Совпадение объясняется тем, что обе сети опираются на одни и те же наиболее информативные признаки.

\begin{table}[H]\centering\small
\caption{Байесовские сети: точность и AUC на тесте}\label{tab:bayesnet}
\begin{tabular}{lcc}\toprule
Модель & Acc test & Доп.\\\midrule
Наш экспертный DAG & 0.669 & --\\
Выученный DAG (HC+BIC) & 0.669 & ребра $Z\to D$, $\textit{online}\to D$, $D\to\textit{size}$, $D\to\textit{age}$\\
Байес-сеть (наш DAG) & -- & AUC test 0.737\\\bottomrule
\end{tabular}
\end{table}

\textit{Повышенная сложность (ROC байес-сети).} AUC байесовской сети составил $0.737$, что сопоставимо с градиентным бустингом и CatBoost и лишь немного уступает логистической регрессии (рис.~\ref{fig:roc_bn}).

\begin{figure}[H]\centering
\includegraphics[width=0.66\textwidth]{figures/roc_with_bn.png}
\caption{ROC-кривые классификаторов и байесовской сети (тест)}\label{fig:roc_bn}
\end{figure}

\zadanie{3.9}{Выберите лучший и худший классификаторы, сформулировав критерий.}
\otvet
В качестве критерия выбран AUC на тестовой выборке как пороговонезависимая мера качества ранжирования, наиболее уместная в нашей задаче. Лучшая модель это логистическая регрессия (AUC $0.744$), худшая случайный лес (AUC $0.721$). Преимущество простой логистической модели объясняется тем, что зависимость склонности к подключению от признаков после преобразований близка к аддитивной, а сильный нелинейный источник вариации ($U$) ненаблюдаем, поэтому гибкие модели не имеют преимущества и более склонны к переобучению.

\zadanie{3.10}{Повышенная сложность: добавьте метод вне scikit-learn.}
\otvet
В анализ включен CatBoost (градиентный бустинг на симметричных деревьях с упорядоченным целевым кодированием). После настройки он дает точность на тесте $0.669$ и AUC $0.739$, практически совпадая с лучшими методами и не превосходя логистическую регрессию. Это подтверждает, что при ненаблюдаемом ключевом факторе сложные модели не дают выигрыша.

\zadanie{3.11}{Дополнительный анализ.}
\otvet
Сводный вывод раздела: предсказуемость подключения ограничена сверху из-за ненаблюдаемого качества менеджмента; различия между методами невелики, а простые и хорошо калиброванные модели предпочтительнее. Это важно для раздела 5, где вероятности подключения используются как вспомогательные оценки склонности.

\section{Регрессия}

\zadanie{4.1}{Отберите признаки для прогнозирования целевой переменной (без переменной воздействия).}
\otvet
Для прогноза выручки $Y$ используются контрольные переменные $\{\textit{age},\textit{size},\textit{online},\textit{city}\}$. Переменная воздействия $D$ исключена согласно заданию, ненаблюдаемое $U$ недоступно. Использованы четыре метода: метод наименьших квадратов (МНК), метод ближайших соседей, случайный лес и градиентный бустинг.

\zadanie{4.2}{Для каждого метода оцените RMSE и MAPE на обучении, тесте и по кросс-валидации; подберите гиперпараметры; сведите в таблицу; обсудите переобучение.}
\otvet
Результаты в таблице~\ref{tab:reg_tuning}. После настройки все методы дают близкое качество (RMSE около $42.6$--$43.5$, MAPE около $0.19$). Случайный лес с параметрами по умолчанию сильно переобучен (RMSE на обучении $18.3$ против $49.8$ на тесте); ограничение глубины и размера листа снижает тестовый RMSE до $43.6$. У МНК и градиентного бустинга разрыв между обучением и тестом мал. Остаточная ошибка велика и после настройки, поскольку существенная часть вариации выручки определяется ненаблюдаемым качеством менеджмента $U$ и случайным шумом: это неустранимая (ирредуцируемая) ошибка, а не дефект моделей.

\begin{table}[H]\centering\small
\caption{RMSE и MAPE регрессионных моделей до и после настройки}\label{tab:reg_tuning}
\begin{tabular}{lccccc}\toprule
Метод & RMSE train & RMSE CV & RMSE test & MAPE test & RMSE test (тюн.)\\\midrule
МНК (OLS) & 42.18 & 42.23 & 42.60 & 0.190 & 42.60\\
kNN & 37.38 & 45.77 & 46.56 & 0.208 & 43.12\\
RandomForest & 18.31 & 49.35 & 49.75 & 0.219 & 43.55\\
GradBoost & 41.01 & 42.33 & 42.67 & 0.190 & 42.72\\\bottomrule
\end{tabular}
\end{table}

\textit{Повышенная сложность (OOB для бустинга).} При обучении бустинга со стохастическим подвыбором ($\textit{subsample}=0.8$) число деревьев подобрано по out-of-bag улучшению функции потерь: оптимум достигается при $472$ деревьях и дает RMSE на тесте $42.86$ против $42.72$ при кросс-валидационной настройке. OOB-критерий вычисляется почти бесплатно в ходе обучения, но опирается на стохастический подвыбор и оценивает только число итераций, тогда как кросс-валидация более универсальна и здесь дала чуть лучший результат.

\zadanie{4.3}{Выберите лучшую и худшую модели и обоснуйте выбор.}
\otvet
По RMSE и MAPE на тесте лучшие результаты у МНК ($42.60$; $0.190$) и градиентного бустинга ($42.72$; $0.189$), худший у случайного леса ($43.55$). Конкурентоспособность МНК объясняется тем, что после преобразований ($\sqrt{\textit{size}}$, $\ln\textit{age}$) базовая выручка почти аддитивна по контролям, а основной нелинейный фактор ($U$) ненаблюдаем, поэтому гибкие модели не имеют преимущества. В качестве лучшей ML-модели для дальнейшего использования в разделе 5 выбран градиентный бустинг (близок к МНК и гибче учитывает возможные нелинейности), в качестве слабой модели метод ближайших соседей.

\zadanie{4.4}{Повышенная сложность: добавьте метод вне scikit-learn.}
\otvet
CatBoost после настройки дает RMSE на тесте $43.33$ и MAPE $0.192$, не превосходя МНК и градиентный бустинг, что согласуется с выводом об ограниченной роли сложных моделей при ненаблюдаемом ключевом факторе.

\zadanie{4.5}{Дополнительный анализ.}
\otvet
Сопоставление разделов 3 и 4 показывает устойчивую картину: при наличии сильного ненаблюдаемого фактора простые, хорошо специфицированные модели не уступают сложным. Это мотивирует акцент раздела 5 не на качестве прогноза, а на корректной идентификации причинного эффекта.

\section{Эффекты воздействия}

\zadanie{5.1}{Запишите выражения для ATE и LATE и проинтерпретируйте их.}
\otvet
Средний эффект воздействия $\mathrm{ATE}=\mathbb{E}[Y_i(1)-Y_i(0)]=\mathbb{E}[\tau_i]$ это ожидаемое изменение месячной выручки при подключении к маркетплейсу для случайно выбранной фирмы. Локальный средний эффект $\mathrm{LATE}=\mathbb{E}[Y_i(1)-Y_i(0)\mid D_i(1)>D_i(0)]$ это средний эффект только для фирм-комплаеров, то есть тех, кто подключается при получении приглашения и не подключается без него. В наших обозначениях $\mathrm{LATE}=\mathbb{E}[\tau_i\mid D_i(0)=0,\,D_i(1)=1]$. Содержательно ATE отвечает на вопрос о пользе подключения для типичной фирмы, а LATE о пользе для тех фирм, чье поведение можно изменить именно приглашением, что особенно важно при оценке программ стимулирования.

\zadanie{5.2}{По симулированным потенциальным исходам получите истинные ATE, CATE и LATE; результаты для ATE и LATE в таблице, для CATE гистограмма.}
\otvet
На супервыборке в $1{,}000{,}000$ наблюдений истинный $\mathrm{ATE}=27.96$ тыс. руб., истинный $\mathrm{LATE}=27.44$ тыс. руб., доля комплаеров $0.264$. Распределение индивидуальных эффектов $\tau_i$ показано на рис.~\ref{fig:cate_true}: эффект существенно неоднороден (примерно от $0$ до $60$ тыс. руб.). Истинные условные эффекты по подгруппам убывают с наличием онлайн-канала и расположением в крупном городе: от $35.99$ для фирм без онлайн-канала вне крупных городов до $19.39$ для фирм с онлайн-каналом в крупных городах, что соответствует заложенным знакам $\tau_2<0$, $\tau_3<0$. Близость ATE и LATE объясняется тем, что отбор в комплаеры в нашей модели слабо смещает распределение $\tau$.

\begin{figure}[H]\centering
\includegraphics[width=0.66\textwidth]{figures/cate_true.png}
\caption{Истинное распределение индивидуальных эффектов воздействия (CATE)}\label{fig:cate_true}
\end{figure}

\zadanie{5.3}{Оцените ATE как разницу средних и опишите недостатки этого подхода.}
\otvet
Наивная разница средних выручки подключившихся и не подключившихся фирм равна $62.91$ тыс. руб. при истинном ATE $27.96$, то есть завышает эффект более чем вдвое. Причина в отборе: к маркетплейсу чаще подключаются фирмы с лучшим менеджментом ($U$), у которых выше и базовая выручка, и сам эффект, поэтому разница средних приписывает подключению часть различий, обусловленных ненаблюдаемыми способностями собственника.

\zadanie{5.4}{Оцените ATE методами OLS, условных мат. ожиданий, IPW, двойной устойчивости и DML; назовите ключевую предпосылку и обсудите ее.}
\otvet
Результаты в таблице~\ref{tab:ate}. Все методы, опирающиеся на корректировку по наблюдаемым контролям, дают близкие оценки около $49$ тыс. руб., снижая смещение наивной разницы ($62.9$) примерно до $+21$, но не устраняя его. Ключевая предпосылка всех этих методов это игнорируемость воздействия (unconfoundedness): при условии на наблюдаемые признаки назначение воздействия не связано с потенциальными исходами. В нашем случае она нарушается, поскольку ключевой фактор отбора, качество менеджмента $U$, ненаблюдаем; поэтому даже двойное машинное обучение, устойчивое к ошибкам спецификации вспомогательных моделей, остается смещенным. Это и есть основной содержательный урок: при наличии ненаблюдаемого конфаундера селекция на наблюдаемых не дает причинной оценки, и интерпретировать $49$ как эффект подключения нельзя.

\begin{table}[H]\centering\small
\caption{Оценки ATE разными методами (вся выборка, $n=10000$)}\label{tab:ate}
\begin{tabular}{lcc}\toprule
Метод оценки ATE & Оценка & Смещение от истинного ATE\\\midrule
Наивная разница средних & 62.91 & $+34.95$\\
OLS (линейная корректировка) & 49.10 & $+21.14$\\
Условные мат. ожидания (g-computation) & 49.07 & $+21.11$\\
IPW (взвешивание на склонность) & 49.40 & $+21.44$\\
AIPW (двойная устойчивость) & 49.42 & $+21.46$\\
Двойное машинное обучение (DML, без IV) & 49.26 & $+21.30$\\
\textbf{Истинный ATE (супервыборка)} & \textbf{27.96} & $0.00$\\\bottomrule
\end{tabular}
\end{table}

\textit{Повышенная сложность.} В качестве дополнительного метода применена двойная устойчивость (AIPW): она остается состоятельной, если верно специфицирована хотя бы одна из двух вспомогательных моделей (исходов или склонности). Здесь это не спасает от смещения, так как нарушена не спецификация, а сама предпосылка игнорируемости.

\zadanie{5.5}{Оцените CATE методами OLS, S-learner, T-learner, трансформации классов и X-learner; сравните и обсудите применимость X-learner и использование оценок.}
\otvet
Сводка в таблице~\ref{tab:cate}; сопоставление оценок X-learner с истинными по подгруппам в таблице~\ref{tab:cate_groups}; распределение оценок на рис.~\ref{fig:cate_est}. Все метаобучатели наследуют смещение уровня из-за ненаблюдаемого конфаундера (средняя оценка около $49$ против истинных $28$), однако сохраняют структуру неоднородности: упорядочение подгрупп по наличию онлайн-канала и расположению воспроизводится верно, а ранговая корреляция оценок с истинными $\tau_i$ положительна и наиболее высока у S-learner и X-learner (около $0.44$). Метод трансформации классов дает самые шумные оценки (ранговая корреляция $0.13$, большой RMSE), что типично для него при умеренных размерах выборки из-за деления на склонность. Применение X-learner оправдано при заметной несбалансированности групп и неоднородности эффекта; в нашем случае группы сбалансированы, поэтому он не дает решающего преимущества перед S-learner. Полученные оценки CATE можно использовать для приоритизации: маркетплейсу выгоднее в первую очередь привлекать фирмы без онлайн-канала вне крупных городов (наибольший эффект), а государственные программы поддержки логично нацеливать на те же сегменты; при этом из-за смещения уровня оценки следует использовать для ранжирования сегментов, а не как абсолютную величину выгоды.

\begin{table}[H]\centering\small
\caption{Оценки CATE: среднее, точность восстановления и ранговая корреляция с истинным эффектом}\label{tab:cate}
\begin{tabular}{lccc}\toprule
Метод CATE & Среднее & RMSE к истинному CATE & Ранг. корр. (Spearman)\\\midrule
OLS-взаимодействия & 49.12 & 24.00 & 0.437\\
S-learner & 49.07 & 23.87 & 0.448\\
T-learner & 49.69 & 24.85 & 0.401\\
Трансформация классов & 49.40 & 50.46 & 0.133\\
X-learner & 49.43 & 24.28 & 0.440\\
\textbf{Истинный CATE} & \textbf{28.20} & 0.00 & 1.000\\\bottomrule
\end{tabular}
\end{table}

\begin{table}[H]\centering
\caption{Истинный и оценённый (X-learner) CATE по подгруппам, тыс. руб.}\label{tab:cate_groups}
\begin{tabular}{llcc}\toprule
online & city & Истинный & X-learner\\\midrule
0 & 0 & 36.01 & 57.58\\
0 & 1 & 26.08 & 47.08\\
1 & 0 & 29.43 & 50.85\\
1 & 1 & 19.63 & 40.50\\\bottomrule
\end{tabular}
\end{table}

\begin{figure}[H]\centering
\includegraphics[width=0.66\textwidth]{figures/cate_est.png}
\caption{Оценённое (X-learner) и истинное распределение CATE}\label{fig:cate_est}
\end{figure}

\zadanie{5.6}{Оцените LATE двойным машинным обучением без инструмента и с инструментом; сопоставьте и проинтерпретируйте.}
\otvet
Результаты в таблице~\ref{tab:late}. Двойное машинное обучение без инструмента опирается на игнорируемость и дает смещенную оценку около $49$, как и в задании 5.4. С инструментом (модель IIVM) оценка LATE равна $25.23$ при истинном $27.44$, то есть смещение сокращается до $-2.2$. Простая оценка Wald без ковариат дает $20.93$ и более изменчива, поскольку использует только инструмент и имеет высокую дисперсию при умеренной доле комплаеров. Таким образом, именно инструментальная переменная, а не корректировка по наблюдаемым, восстанавливает причинный эффект. В нашем случае ATE и LATE близки по величине, поэтому основное различие между оценками без инструмента и с инструментом объясняется не разницей между средним и локальным эффектами, а устранением смещения от ненаблюдаемого конфаундера. Содержательно $\mathrm{LATE}\approx 25$--$27$ тыс. руб. это прирост выручки для фирм, которых к подключению подталкивает именно приглашение; на них и стоит ориентировать программы стимулирования.

\begin{table}[H]\centering\small
\caption{Оценки LATE (вся выборка, $n=10000$)}\label{tab:late}
\begin{tabular}{lcc}\toprule
Метод оценки LATE & Оценка & Смещение от истинного LATE\\\midrule
Wald / 2SLS (без ковариат) & 20.93 & $-6.51$\\
DML без IV (смещён, не LATE) & 49.26 & $+21.82$\\
DML с IV & 25.23 & $-2.21$\\
\textbf{Истинный LATE (супервыборка)} & \textbf{27.44} & $0.00$\\\bottomrule
\end{tabular}
\end{table}

\textit{Повышенная сложность (параметрическая модель).} В качестве параметрической альтернативы мы оценили модель эндогенного переключения (switching regression) методом максимального правдоподобия. Она состоит из уравнения отбора (подключения) с инструментом $Z$ и двух уравнений выручки для режимов $D=0$ и $D=1$, а совместная нормальность ошибок с корреляциями $\rho_0,\rho_1$ между уравнением отбора и уравнениями исходов корректирует эндогенность. Оцененные корреляции значимо отличны от нуля ($\hat\rho_1=0.67$, $\hat\rho_0=-0.47$), что подтверждает наличие эндогенного отбора. Модель дает оценку среднего эффекта $35.3$ тыс. руб. (таблица~\ref{tab:switch}): она устраняет существенную часть смещения по сравнению с методами на наблюдаемых ($49$), но остается смещенной вверх (на $+7.4$ относительно истинного ATE), поскольку предпосылки нормальности и линейности не соответствуют нелинейному процессу генерации данных (квадратичный вклад $U$ и зависимость эффекта от $U$). Это иллюстрирует ключевой компромисс: параметрический подход эффективен при верной спецификации совместного распределения, но уязвим к ее нарушению, тогда как полупараметрическое двойное машинное обучение с инструментом устойчивее и здесь оказалось точнее. Тот же класс моделей реализован, в частности, в пакете \texttt{switchSelection} (Коссова, Потанин, 2018, 2022); в нашей работе он полностью реализован средствами Python, на котором, в соответствии с требованиями, выполнены все задания.

\begin{table}[H]\centering\small
\caption{Параметрическая модель эндогенного переключения (MLE) в сравнении}\label{tab:switch}
\begin{tabular}{lcc}\toprule
Метод & Оценка эффекта & Смещение от истинного ATE\\\midrule
Методы на наблюдаемых (DML без IV) & 49.26 & $+21.30$\\
Параметрическое эндог. переключение (MLE) & 35.32 & $+7.36$\\
Двойное машинное обучение с IV & 25.23 & $-2.73$\\
\textbf{Истинный ATE} & \textbf{27.96} & $0.00$\\\bottomrule
\end{tabular}
\end{table}

\zadanie{5.7}{Оцените ATE и LATE на худших моделях и сравните с лучшими; сделайте вывод об устойчивости.}
\otvet
В таблице~\ref{tab:robust} вспомогательные модели двойного машинного обучения заменены со связки лучших (градиентный бустинг для исходов и логистическая регрессия для склонности) на связку слабых (метод ближайших соседей и случайный лес). Оценки меняются незначительно: ATE остается около $49$, LATE около $25$--$26$. Это согласуется с теорией двойного машинного обучения: благодаря ортогонализации и перекрестной подгонке итоговая оценка слабо чувствительна к умеренному ухудшению вспомогательных моделей. Важный вывод: устойчивость к качеству ML-моделей не спасает от смещения, вызванного нарушением идентифицирующих предпосылок, и именно корректная идентификация (через инструмент), а не выбор алгоритма, определяет достоверность причинной оценки.

\begin{table}[H]\centering\small
\caption{Устойчивость оценок к качеству вспомогательных ML-моделей}\label{tab:robust}
\begin{tabular}{lcccc}\toprule
Вспомогательные модели & DML без IV (ATE) & Смещение ATE & DML с IV (LATE) & Смещение LATE\\\midrule
Лучшие (бустинг + логит) & 49.26 & $+21.30$ & 25.23 & $-2.21$\\
Слабые (kNN + случайный лес) & 48.69 & $+20.73$ & 26.28 & $-1.16$\\\bottomrule
\end{tabular}
\end{table}

\zadanie{5.8}{Резюмируйте ключевые выводы раздела.}
\otvet
Во-первых, наивная разница средних и любые методы на наблюдаемых признаках (OLS, условные мат. ожидания, IPW, двойная устойчивость, DML без инструмента) сильно завышают эффект подключения ($\approx 49$ против истинного $\approx 28$) из-за ненаблюдаемого качества менеджмента. Во-вторых, корректную оценку дает инструментальная переменная: двойное машинное обучение с инструментом восстанавливает LATE ($25.2$ при истинном $27.4$). В-третьих, метаобучатели не дают несмещенной величины эффекта, но верно ранжируют сегменты по силе эффекта, что полезно для таргетинга. В-четвертых, результаты устойчивы к качеству вспомогательных ML-моделей, но не к нарушению идентифицирующих предпосылок. Практический вывод: для оценки отдачи от подключения к маркетплейсу решающую роль играет наличие источника экзогенной вариации (случайные приглашения), а не сложность алгоритма.

\zadanie{5.9}{Дополнительный анализ.}
\otvet
Дополнительно отметим, что близость ATE и LATE в нашем случае не является общим правилом: при более сильной зависимости индивидуального эффекта от факторов, определяющих комплаентность, эти величины расходятся, и тогда выбор между ними становится содержательно значимым для политики стимулирования.

\newpage
\section*{Список литературы}
\addcontentsline{toc}{section}{Список литературы}

\refitem{Долгих, С. И., \& Потанин, Б. С. (2023). Влияние государственного управления на эффективность российских фирм. \textit{Проблемы прогнозирования}, (1), 90--103.}
\refitem{Коссова, Е. В., \& Косорукова, М. А. (2023). Оценивание влияния высшего образования на здоровье: сравнение многомерной рекурсивной пробит-модели и мэтчинга. \textit{Прикладная эконометрика, 69}, 65--90.}
\refitem{Коссова, Е. В., \& Потанин, Б. С. (2018). Обобщение метода Хекмана и модели с переключением на случай произвольного числа уравнений отбора. \textit{Прикладная эконометрика, 50}, 114--143.}
\refitem{Коссова, Е. В., \& Потанин, Б. С. (2022). Оценивание модели с гауссовским мультиномиальным эндогенным переключением и неслучайным отбором. \textit{Прикладная эконометрика, 67}(3), 121--143.}
\refitem{Котырло, Е. С. (2024). Простой и сложный метод разности разностей. \textit{Прикладная эконометрика, 73}(1), 119--142.}
\refitem{Couture, V., Faber, B., Gu, Y., \& Liu, L. (2021). Connecting the countryside via e-commerce: Evidence from China. \textit{American Economic Review: Insights, 3}(1), 35--50. \url{https://doi.org/10.1257/aeri.20190382}}
\refitem{Imbens, G. W., \& Angrist, J. D. (1994). Identification and estimation of local average treatment effects. \textit{Econometrica, 62}(2), 467--475. \url{https://doi.org/10.2307/2951620}}
\refitem{Lendle, A., Olarreaga, M., Schropp, S., \& Vézina, P.-L. (2016). There goes gravity: eBay and the death of distance. \textit{The Economic Journal, 126}(591), 406--441. \url{https://doi.org/10.1111/ecoj.12286}}

\end{document}
